
\subsection*{2.8. Unipotent group schemes over a field}
\phantomsection
\label{sga3:VIIB:2.8}
\addcontentsline{toc}{subsection}{\SGARefLink{sga3:VIIB:2.8}{2.8}. Unipotent group schemes over a field}
\footnote{\label{sga3:VIIB:nde:115}The editors added this subsection.}

Let \(k\) be a field. Recall that a \(k\)-group scheme \(G\) is
called \emph{unipotent} if it satisfies the following two conditions
\textup{(cf.\ [DG70], \S~IV.2, \SGARefLink{sga3:VIIB:2.5:proposition}{Proposition 2.5})}:
\begin{enumerate}
\item[\textup{(a)}]
\(G\) is affine.

\item[\textup{(b)}]
Every simple \(\mathcal O(G)\)-comodule is trivial. Explicitly, if
\[
  \rho:V\longrightarrow V\otimes_k\mathcal O(G)
\]
is an \(\mathcal O(G)\)-comodule structure on a nonzero
\(k\)-vector space \(V\), and there is no nonzero proper subspace
\(W\subset V\) such that
\[
  \rho(W)\subset W\otimes_k\mathcal O(G),
\]
then \(V\) is one-dimensional and
\[
  \rho(v)=v\otimes1
  \qquad\text{for every }v\in V.
\]
\end{enumerate}
According to the cited result, when \(G\) is of finite type over
\(k\), this is equivalent to the definition given in
Exposé XVII, \S~1: \(G\) has a finite composition series whose
successive quotients are isomorphic to \(k\)-subgroups of
\(\mathbb G_{a,k}\).

For every affine \(k\)-group scheme \(G\), the comultiplication of
\(\mathcal O(G)\) endows the pseudocompact \(k\)-module
\[
  A=\mathcal O(G)^*
\]
with the structure of a profinite \(k\)-algebra, not necessarily
commutative; its unit \(1_A\) is the augmentation
\[
  \varepsilon:\mathcal O(G)\longrightarrow k.
\]
Set
\[
  I=\{f\in A\mid f(1_{\mathcal O(G)})=0\}.
\]
Then \(I\) is a two-sided ideal of \(A\), and
\[
  A=k\cdot1_A\oplus I;
\]
cf.\ \SGARefLink{sga3:VIIB:1.3.6}{1.3.6}.

Let \(V\) be a finite-codimensional subspace of \(A\), and consider
the continuous \(k\)-bilinear map
\[
  \varphi:A\times A\longrightarrow A/V,
  \qquad
  (a,b)\longmapsto ab+V.
\]
By \SGARefLink{sga3:VIIB:0.3.1}{Lemma 0.3.1}, there are finite-codimensional subspaces
\(L_1,L_2\subset A\) such that \(V\) contains \(AL_2\) and
\(L_1A\). Thus
\[
  L=L_1\cap L_2
\]
is finite-codimensional, and \(V\) contains the two-sided ideal
\(ALA\), which is itself finite-codimensional. Hence the
finite-codimensional two-sided ideals form a neighborhood basis of
zero in \(A\).

It follows that an \(\mathcal O(G)\)-comodule is the same thing as
a continuous \(A\)-module: an \(A\)-module \(V\), equipped with the
discrete topology, for which the map
\[
  A\times V\longrightarrow V
\]
is continuous. Every such module is the union of
finite-dimensional \(A\)-submodules \(V_i\), each of which is a
module over a finite-dimensional quotient \(k\)-algebra \(A_i\) of
\(A\).

Consequently, a simple continuous \(A\)-module \(M\) is
finite-dimensional over \(k\) and is a faithful simple module over
the finite-dimensional algebra \(A/\operatorname{Ann}(M)\). The
latter is therefore a finite-dimensional simple \(k\)-algebra, so
\(\operatorname{Ann}(M)\) is an open maximal ideal of \(A\).

Conversely, let \(P\) be an open prime ideal of
\(A\).\footnote{\label{sga3:VIIB:nde:116}Recall that a two-sided
ideal \(P\) of a ring \(A\) is called prime if, in \(A/P\), the
product of any two nonzero two-sided ideals is nonzero.}
Then \(A/P\) is a finite-dimensional \(k\)-algebra in which the
zero ideal is prime, and hence is a finite-dimensional simple
\(k\)-algebra. Up to isomorphism, there is therefore a unique
simple continuous \(A\)-module with annihilator \(P\). Thus
\[
  M\longmapsto\operatorname{Ann}(M)
\]
is a bijection from the isomorphism classes of simple continuous
\(A\)-modules to the open prime ideals of \(A\).

In particular, the \(A\)-module \(A/I\), which is one-dimensional
over \(k\), is called the \emph{trivial module}. It corresponds to
the trivial one-dimensional \(\mathcal O(G)\)-comodule \(V\), for
which
\[
  \rho(v)=v\otimes1_{\mathcal O(G)}
  \qquad\text{for every }v\in V.
\]
We have proved:

\medskip
\noindent\textbf{Proposition 2.8.1.}
\phantomsection
\label{sga3:VIIB:2.8.1}
{\itshape Let \(k\) be a field, let \(G\) be an affine
\(k\)-group scheme, and put \(A=\mathcal O(G)^*\).
\begin{enumerate}
\item[\textup{(i)}]
\(G\) is unipotent if and only if \(I\) is the unique open prime
ideal of \(A\).\footnote{\label{sga3:VIIB:nde:117}Equivalently,
\(k\cdot1_{\mathcal O(G)}\) is the unique simple
\(k\)-subcoalgebra of \(\mathcal O(G)\); see, for example,
[Ab80], 3.1.4. The editors also note a misprint in the cited work,
p.~130, line~4: \(M\simeq C^*/\operatorname{ann}M\) should be
replaced by
\(C^*/\operatorname{ann}M\simeq\operatorname{End}_k(M)\).}

\item[\textup{(ii)}]
In particular, suppose that \(G\) is commutative, so that
\[
  \mathcal O(G)=\mathbf H(D(G)),
\]
where \(D(G)=\operatorname{Spf}(A)\) is the Cartier dual of \(G\).
Then \(G\) is unipotent if and only if \(D(G)\) is infinitesimal.
\end{enumerate}}

